\documentclass[a4paper]{article}     

\usepackage{amsfonts}
\usepackage{amsmath}
\usepackage{amssymb}
\usepackage{graphicx}
\usepackage{color}

\newcommand{\Q}{\mathbb{Q}}
\newcommand{\R}{\mathbb{R}}
\newcommand{\llim}{\lim\limits}
\newcommand{\dint}{\displaystyle\int}

\begin{document}

\noindent \large Auteur: Sadia Nazary \\
\noindent \large Studierichting: Informatica\\
\noindent \large Oefeningenreeks 6D nr. 31.15\\

\medskip

\normalsize

$\displaystyle\sum\limits_{n=2}^{\infty}\dfrac{1}{(\ln n)^{\ln n}}$\\

We schrijven de algemene term anders:\\

$\dfrac{1}{(\ln n)^{\ln n}}
=
e^{-\ln n\ln(\ln n)}$\\

$=
\dfrac{1}{n^{\ln(\ln n)}}$\\

Vanaf $n>e^{e^2}$ geldt:\\

$\ln(\ln n)>2$\\

Dus:\\

$\dfrac{1}{n^{\ln(\ln n)}}<\dfrac{1}{n^2}$\\

Aangezien $\displaystyle\sum\limits_{n=2}^{\infty}\dfrac{1}{n^2}$ convergeert, convergeert ook\\

$\displaystyle\sum\limits_{n=2}^{\infty}\dfrac{1}{(\ln n)^{\ln n}}$\\

door het vergelijkingscriterium.\\

$\boxed{\text{Convergent}}$

\begin{thebibliography}{999}
%\bibitem{a} Referentie
\end{thebibliography}

\end{document}